% !TeX root = main.tex
\begin{thebibliography}{11}
\addcontentsline{toc}{section}{Referencias}
\small

\bibitem{szeliski2011}
R. Szeliski (2011).
\textit{Computer Vision: Algorithms and Applications}.
Primera edición, Springer.
\href{https://doi.org/10.1007/978-1-84882-935-0}{doi:10.1007/978-1-84882-935-0}.

\bibitem{boykov2001}
Y. Y. Boykov y M.-P. Jolly (2001).
Interactive Graph Cuts for Optimal Boundary \& Region Segmentation
of Objects in N-D Images.
En \textit{Proceedings of the Eighth IEEE International Conference on Computer Vision}, vol.~1, pp.~105--112.
\href{https://doi.org/10.1109/ICCV.2001.937505}{doi:10.1109/ICCV.2001.937505}.

\bibitem{rother2004}
C. Rother, V. Kolmogorov y A. Blake (2004).
``GrabCut'': Interactive Foreground Extraction Using Iterated Graph Cuts. \textit{ACM Transactions on Graphics}, 23(3), 309--314.
\href{https://doi.org/10.1145/1015706.1015720}{doi:10.1145/1015706.1015720}.

\bibitem{shapeydatos}
J. Shapey, A. Kujawa, R. Dorent, G. Wang, S. Bisdas, A. Dimitriadis,
D. Grishchuck, I. Paddick, N. Kitchen, R. Bradford, S. Saeed, S. Ourselin
y T. Vercauteren (2021).
Segmentation of Vestibular Schwannoma from Magnetic Resonance Imaging:
An Open Annotated Dataset and Baseline Algorithm.
Versión~2, conjunto de datos \textit{Vestibular-Schwannoma-SEG},
The Cancer Imaging Archive, licencia CC BY 4.0.
\href{https://doi.org/10.7937/TCIA.9YTJ-5Q73}{doi:10.7937/TCIA.9YTJ-5Q73}.

\bibitem{clark2013}
K. Clark, B. Vendt, K. Smith, J. Freymann, J. Kirby, P. Koppel, S. Moore,
S. Phillips, D. Maffitt, M. Pringle, L. Tarbox y F. Prior (2013).
The Cancer Imaging Archive (TCIA): Maintaining and Operating a Public
Information Repository.
\textit{Journal of Digital Imaging}, 26(6), 1045--1057.
\href{https://doi.org/10.1007/s10278-013-9622-7}{doi:10.1007/s10278-013-9622-7}.

\bibitem{shapey2021}
J. Shapey, A. Kujawa, R. Dorent, G. Wang, A. Dimitriadis, D. Grishchuk,
I. Paddick, N. Kitchen, R. Bradford, S. R. Saeed, S. Bisdas, S. Ourselin
y T. Vercauteren (2021).
Segmentation of vestibular schwannoma from MRI, an open annotated
dataset and baseline algorithm.
\textit{Scientific Data}, 8, artículo~286.
\href{https://doi.org/10.1038/s41597-021-01064-w}{doi:10.1038/s41597-021-01064-w}.

\bibitem{presentacion}
F. Castro Chaves y E. Chacón Araya (2026).
\textit{Segmentación interactiva de tejidos con GrabCut: schwannoma
vestibular en resonancia magnética}.
Presentación del Foro Tema~4, PF-5039 Análisis de Imágenes, IIS-2026,
diapositivas~10--12 y~16. Material del proyecto no publicado.

\bibitem{cuaderno}
Proyecto \textit{horus-grabcut} (2026).
Cuaderno de segmentación interactiva \path{notebooks/grabcut.ipynb}
y archivos de \path{data/}.
Copia local consultada el 10 de octubre de 2026; se utilizan sus salidas
guardadas para \texttt{VS-SEG-001}. La copia local del cuaderno contiene
modificaciones respecto de la versión registrada en Git.
\href{https://github.com/DarthFranco2001/horus-grabcut}{Repositorio del proyecto}.

\bibitem{dinic1970}
E. A. Dinic (1970).
Algorithm for Solution of a Problem of Maximum Flow in Networks
with Power Estimation. \textit{Soviet Mathematics Doklady}, 11, 1277--1280.
\href{https://www.cs.bgu.ac.il/~dinitz/D70.pdf}{Texto de la publicación en el sitio del autor}.

\bibitem{repositorio}
Proyecto \textit{horus-grabcut} (2026).
\textit{horus-grabcut}. Repositorio de código fuente: cuaderno,
aplicación web, datos de trabajo y pruebas.
\url{https://github.com/DarthFranco2001/horus-grabcut}.

\bibitem{web}
Proyecto \textit{horus-grabcut} (2026).
\textit{Laboratorio GrabCut}. Aplicación web interactiva.
Consultada el 10 de octubre de 2026.
\url{https://darthfranco2001.github.io/horus-grabcut/}.

\end{thebibliography}
